\section{总结与延伸}

这堂课把 RL 从优化算法重新放回产品系统。Product belief 决定希望看到的行为；interface 暴露任务与失败；eval 把“感觉”拆成可复现 evidence；data debugging 追踪行为来源；RL environment 定义状态、动作、工具与恢复；reward 把结果和过程写进目标；上线后的用户 evidence 又启动下一轮研究。任何一层缺失，都会让另一层承担它无法解决的问题。

\subsection{一张表串起 co-design loop}

\begin{table}[H]
\centering
\small
\begin{tabular}{@{}L{0.16\textwidth}L{0.24\textwidth}L{0.27\textwidth}L{0.23\textwidth}@{}}
\toprule
层次 & 核心对象 & 代表性失败 & 验收问题 \\
\midrule
Belief & 目标行为与用户关系 & 口号不可操作 & 是否能写成行为 contract \\
Interface & artifact、diff、控制与权限 & 能力不可见或不可撤销 & 用户能否理解、编辑、停止 \\
Eval & 分布、rubric、judge、阈值 & benchmark 与真实任务错位 & 是否可复现、分 slice、可决策 \\
Data & provenance、mixture、preference & 数据冲突与污染 & 能否定位行为来源并回滚 \\
Environment & 状态、动作、工具、终止 & 训练任务与部署任务不同 & 权限、预算、失败恢复是否一致 \\
Reward & outcome、process、cost、risk & reward hacking / Goodhart & verifier 是否独立且抗攻击 \\
Deployment & telemetry、回归、责任人 & demo 成功掩盖长尾事故 & 是否监控 outcome 并可快速回滚 \\
\bottomrule
\end{tabular}
\caption{RL、产品与研究协同设计的七层验收地图。}
\end{table}

\subsection{开始一个 Frontier Product 项目前应回答}

\begin{importantbox}{十二问自测}
\begin{enumerate}
  \item Product belief 能否写成可观察行为，而不是形容词？
  \item 目标用户的真实任务和失败成本是什么？
  \item 模型 capability 通过什么 form factor 暴露？
  \item 哪些 artifact、状态和来源必须可追踪？
  \item 用户保留哪些确认、编辑、停止和 rollback 权力？
  \item Eval 的分布、judge、rubric 与版本是否固定？
  \item 是否同时报告 overall、slice、worst-case 与严重案例？
  \item 失败 taxonomy 能否映射到数据、系统或 policy 根因？
  \item RL environment 的工具、权限、预算和终止是否与部署一致？
  \item Reward 是否把真实 outcome、过程、成本与风险分开？
  \item Verifier 有哪些 blind spot，策略是否能攻击它？
  \item 上线后谁看 outcome、谁能回滚、什么条件触发停机？
\end{enumerate}
\end{importantbox}

\subsection{复现记录的最小集合}

无论研究 data、behavior、RL environment 还是 interface，都应保存 source/version、prompt/system、模型与 sampling、工具 schema、完整 trace、judge/rubric、人工 disagreement、失败样本和回滚条件。若使用 AI judge 或 synthetic data，还要固定生成模型、judge 模型与权限；否则一次结果无法被下一版本解释。

\begin{lstlisting}[caption={Frontier product 实验的可复现记录}]
experiment = {
    "hypothesis": product_belief_as_testable_behavior,
    "model": [base, post_training, checkpoint, sampling],
    "environment": [tools, permissions, budget, termination],
    "interface": [surface, artifact_schema, confirmation, rollback],
    "evaluation": [datasets, slices, judges, rubrics, thresholds],
    "evidence": [traces, outcomes, disagreements, severe_failures],
    "decision": [ship, hold, rollback, next_hypothesis],
}
\end{lstlisting}

\subsection{拓展阅读}

\begin{itemize}
  \item Bai et al., \emph{Constitutional AI: Harmlessness from AI Feedback}：理解 RLAIF、原则约束与 AI feedback。
  \item Röttger et al., \emph{XSTest}：研究 exaggerated safety 与 benign prompt over-refusal。
  \item Lilian Weng, \emph{Reward Hacking in Reinforcement Learning}：系统整理 proxy reward 与 hacking 路径。
  \item Ouyang et al., \emph{Training language models to follow instructions with human feedback}：RLHF 的经典 pipeline 与局限。
  \item Amodei et al., \emph{Concrete Problems in AI Safety}：从 reward hacking、side effects 与 oversight 理解环境设计。
\end{itemize}

\begin{warningbox}{最后的边界}
本讲中的产品名、模型表现、成本曲线和组织实践都属于 2025 年课堂时间点。讲义保存的是 co-design 方法，不把任何单个 demo、benchmark 或公司流程当作永久事实。真正可迁移的能力是：把愿景翻译成环境，把行为翻译成 eval，把失败追到数据与系统，再用可复现证据决定是否交付。
\end{warningbox}

\end{document}
